\documentclass[11pt]{article}
\usepackage[margin=1in]{geometry}
\usepackage{amsmath,amssymb,amsthm}
\usepackage[round]{natbib}
\usepackage{hyperref}
\usepackage{microtype}

\newtheorem{theorem}{Theorem}
\newtheorem{proposition}{Proposition}
\newtheorem{lemma}{Lemma}
\newtheorem{corollary}{Corollary}
\newtheorem{conjecture}{Conjecture}
\newtheorem{remark}{Remark}

\title{The Width of the Conformal Fan:\\
Dependence and the Variance of Realized Coverage}
\author{Peter Cotton\\\texttt{peter.cotton@microprediction.com}}
\date{June 2026}

\begin{document}
\maketitle

\begin{abstract}
Fix a split-conformal calibration set; the coverage realized against an independent draw from the
score marginal is $c=U_{(k)}$, distributed $\mathrm{Beta}(k,n-k+1)$ for independent calibration
scores, with variance about $\alpha(1-\alpha)/n$ --- the fan. We show the fan's width is set by
the dependence among the calibration scores, through the covariance kernel of the sub-threshold
count. Three results follow: a leading-order fan coefficient in a Bahadur regime; an exact
aggregate bound, summed over all levels, showing negative association never widens the fan; and,
for equicorrelated normals, the exact single-level law
$\operatorname{Var}(Z_{(k)})=v_k+\rho(1-v_k)$. The general single-level contraction is false: an explicit negatively associated law, uniform marginals, over-disperses an order
statistic at every $n$, and full support does not rescue it (a strictly-positive-density,
strictly-interior-NA law violates too, certified exactly). The bound does hold, at every level, on
the designs a practitioner actually builds: sampling without replacement gives the exact closed form
$\operatorname{Var}(U_{(k)})=\tfrac{M-n}{M+1}V_0$ with the coverage mean on target, stratified and
Latin-hypercube designs contract with exact variances, and the FGM and general negative-correlation
Gaussian copulas contract by a monotonicity argument that also closes the transfer to the uniform
scale. Positive dependence --- serial correlation, clustering, shared factors, the common case ---
runs the mirror: the fan reliably \emph{widens}, at every level in every structured family, and we
conjecture it can never narrow a single level, a genuine asymmetry with the negative side.
\end{abstract}

\section{Setup}

Split conformal returns a set with marginal coverage $1-\alpha$. Fix the calibration set; the
coverage you then realize on fresh test points is a random variable, and it is not $1-\alpha$ on the
nose. Write the scores on the probability-integral scale, $U_i=F(S_i)$. We assume throughout that the
score marginal $F$ is continuous, so each $U_i\sim\mathrm{Unif}(0,1)$; if $F$ has atoms, read $U_i$ as
the randomised PIT $F(S_i^-)+V_i\,(F(S_i)-F(S_i^-))$ with $V_i\sim\mathrm{Unif}(0,1)$ independent,
which restores exact uniformity. With $k=\lceil(n+1)(1-\alpha)\rceil$ the conformal threshold is the
order statistic $U_{(k)}$, and the calibration-conditional coverage against an independent test draw
from the marginal is
\begin{equation}
c=U_{(k)} .
\end{equation}
For \emph{independent} calibration scores this is the classical result \citep{vovk2012,marques2024}:
$c\sim\mathrm{Beta}(k,n-k+1)$, with
\begin{equation}\label{eq:betavar}
\mathbb E[c]=\frac{k}{n+1},\qquad \operatorname{Var}(c)=\frac{k(n-k+1)}{(n+1)^2(n+2)}\;\approx\;
\frac{\alpha(1-\alpha)}{n}.
\end{equation}
This across-sample spread of realized coverage is its \emph{dispersion}, the formal object of this
note; the family of such laws swept out as the calibration dependence varies is the \emph{fan}, for
short. The per-dataset coverage
scatters around $1-\alpha$ with variance of order $1/n$ (standard deviation of order $n^{-1/2}$). The
question of this note is what the dependence among the calibration scores does to it.

Two cautions fix the object of study. First, $c=U_{(k)}$ is coverage against a test draw
\emph{independent} of the calibration set; it is not the conformal conditional coverage
$\mathbb P(U_{n+1}\le U_{(k)}\mid U_1,\dots,U_n)$ for a test point exchangeable with, and dependent
on, the calibration scores, which need not equal $U_{(k)}$. The marginal conformal guarantee comes
from the rank symmetry of $U_{n+1}$ among the $n+1$ scores and is not what controls $c$. Second, the
mean $\mathbb E[c]=k/(n+1)$ is an \emph{independence} statement, not a consequence of uniform
marginals or exchangeability: if $U_1=\dots=U_n=V\sim\mathrm{Unif}(0,1)$ then $U_{(k)}=V$ and
$\mathbb E[c]=\tfrac12$ for every $k$. So under dependence both the mean and the variance of $c$
respond; this note studies the variance, and flags the mean where it matters. The companion notes
treat the orthogonal, within-dataset axis \citep{cotton_marginally,cotton_signed}; here every
statement is about the across-dataset fluctuation of $c$. The results separate by type:
Theorems~\ref{thm:pos} and \ref{thm:agg} and Propositions~\ref{prop:path}--\ref{prop:gauss} are
finite-sample exact, Theorem~\ref{thm:coeff} is asymptotic and conditional on a Bahadur
representation, and the general per-level negative-association contraction is \emph{false}
(Proposition~\ref{prop:counter}), holding only on structured subclasses.

A third point fixes what kind of dependence is in scope, since it has to square with an independent
test. What the analysis requires is only that the test draw be independent of the calibration scores;
the calibration dependence is then necessarily \emph{design-induced} --- a property of how the
calibration set is assembled, not of a data stream the test belongs to. This is the standard
split-conformal test assumption (a fresh, independent future point); the only non-standard ingredient
is letting the \emph{calibration} set carry an engineered dependence. That ingredient is not exotic ---
it is what you get whenever the calibration set is \emph{constructed or curated} rather than taken as an
i.i.d.\ draw, and the fan is then a statement about a design choice.

Three concrete settings. \emph{Space-filling calibration.} When the calibration inputs can be chosen
--- active or pool-based calibration, or any pipeline that selects which points to label --- one spreads
them evenly, on a stratified, Latin-hypercube, or low-discrepancy (quasi-Monte-Carlo) grid, to avoid
clumping. Even coverage of the input space makes the resulting scores negatively associated with uniform
marginals, and the fan says this engineering can only tighten the run-to-run variability of realized
coverage, never loosen it. \emph{Balanced subsampling from a pool.} Labeled data is often subsampled to
a fixed calibration budget; drawing without replacement, or balancing on a stratifying covariate, or
antithetically pairing, all induce negative association among the retained scores while the deployed test
point remains exogenous. \emph{Deduplication and diversity filtering.} Removing near-duplicates or
enforcing diversity in a calibration set --- routine when data are scraped or redundant --- is again a
negative-dependence-inducing design, and again a fan-narrowing one. In each case the practical payoff is
the same: a calibration design that is more evenly spread buys a more reliable realized coverage from one
deployment to the next, at no cost to the marginal guarantee, which the rank symmetry of the independent
test point still delivers.

The contrast is the case where the test shares its source with the calibration --- the next step of a
time series, or the unsampled remainder of a without-replacement draw used as the test --- which couples
the test to the calibration, so that $c\ne U_{(k)}$ and a different analysis applies; that case is out of
scope here. The maximally negatively associated floor is read the same way: a deterministic grid laid
down by design, not a draw whose leftover is the test.

The whole problem passes through the count process $N(u)=\sum_{i=1}^n\mathbf 1\{U_i\le u\}$, because
$U_{(k)}=\int_0^1\mathbf 1\{N(u)<k\}\,\mathrm du$, so
\begin{equation}
\operatorname{Var}(c)=\int_0^1\!\!\int_0^1\operatorname{Cov}\big(\mathbf 1\{N(u)<k\},\,
\mathbf 1\{N(v)<k\}\big)\,\mathrm du\,\mathrm dv,
\end{equation}
a functional of the law of the count process alone --- equivalently, by Sklar's theorem, of the score
\emph{copula}, since the uniform marginal has factored out (we return to this in
Remark~\ref{rmk:copula}). Two facts about $N$ drive everything.

\begin{lemma}[Count variance]\label{lem:count}
If the scores are negatively associated \citep{joagdev1983}, then for every $u$,
$\operatorname{Var}(N(u))\le n\,u(1-u)$, with equality iff the sub-threshold indicators are pairwise
uncorrelated.
\end{lemma}
\begin{proof}
The indicators $\mathbf 1\{U_i\le u\}$ are functions of single distinct coordinates, all nondecreasing
in the same direction, so negative association is preserved and they are pairwise non-positively
correlated. So
$\operatorname{Var}(N(u))=\sum_i\operatorname{Var}+\sum_{i\ne j}\operatorname{Cov}\le
\sum_i u(1-u)=n\,u(1-u)$.
\end{proof}

\section{The fan coefficient}

Let $p=1-\alpha$ and $b_i=\mathbf 1\{U_i\le p\}$, so $N(p)=\sum_i b_i$. (We reserve $\xi$ for the
i.i.d.\ normals of Proposition~\ref{prop:gauss}.) The representation is on the uniform scale, where the
quantile density is $1/f(p)=1$, so the usual $1/f$ factor is absent.

\begin{theorem}[Fan coefficient]\label{thm:coeff}
Suppose $F$ is continuous (or a randomised PIT is used), $k/n\to p\in(0,1)$, and the scores satisfy a
Bahadur representation $U_{(k)}-p=p-N(p)/n+r_n$ with the remainder negligible in $L^2$ relative to
$n^{-1}\operatorname{Var}(N(p))^{1/2}$. Then
\begin{equation}
\operatorname{Var}(c)=\frac1{n^2}\operatorname{Var}(N(p))+o\!\Big(\frac{\operatorname{Var}(N(p))}{n^2}\Big).
\end{equation}
If moreover $\operatorname{Var}(N(p))\asymp n$, writing it as
$np(1-p)+n(n-1)\bar c$ with $\bar c=\tfrac1{n(n-1)}\sum_{i\ne j}\operatorname{Cov}(b_i,b_j)$,
\begin{equation}
\operatorname{Var}(c)=\frac{p(1-p)}{n}\,R\,(1+o(1)),\qquad R=1+\frac{n-1}{p(1-p)}\,\bar c .
\end{equation}
\end{theorem}

So the leading fan coefficient is the count variance at the operating quantile. Under
independence $\bar c=0$, $R=1$, recovering \eqref{eq:betavar}; negative association makes $\bar c<0$
and $R<1$, narrowing the fan. Two caveats the statement respects. The multiplicative form needs
$\operatorname{Var}(N(p))\asymp n$: under strong negative dependence the count variance can be
$o(n)$ (zero at the floor), and then only the additive form is safe. And the Bahadur regime itself
fails under a fixed positive latent (de Finetti) dependence, where $N(p)/n$ has a non-degenerate
random limit $Q([0,p])$ and $U_{(k)}$ converges to a random conditional quantile rather than
linearizing at $p$; the next section handles that regime exactly and separately.

\section{Positive dependence: an exact decomposition}

When the scores are extendable, drawn from an infinitely exchangeable sequence, de Finetti
\citep{definetti1937,kerns2006} makes them conditionally independent given a latent directing measure
$Q$. The law of total variance then gives an exact decomposition.

\begin{theorem}[Positive-dependence decomposition]\label{thm:pos}
For an extendable exchangeable calibration sequence with directing measure $Q$,
\begin{equation}
\operatorname{Var}(c)=\underbrace{\mathbb E_Q\!\big[\operatorname{Var}(U_{(k)}\mid Q)\big]}
_{\text{average within-}Q\text{ order-stat variance}}
+\underbrace{\operatorname{Var}_Q\!\big(m_k(Q)\big)}_{\text{between-}Q\text{ term}},
\qquad m_k(Q)=\mathbb E[U_{(k)}\mid Q].
\end{equation}
The between-$Q$ term is non-negative and vanishes under independence.
\end{theorem}

The decomposition is exact and needs no asymptotics, but note what conditioning gives: $U_i\mid Q$ are
i.i.d.\ $Q$, so $U_{(k)}\mid Q$ is the $k$-th order statistic of $Q$ on the marginal scale, \emph{not}
a uniform $\mathrm{Beta}$ variable, and its mean $m_k(Q)$ moves with $Q$. The within-$Q$ term is thus
not the independent fan \eqref{eq:betavar}, and the identity by itself does \emph{not} prove
$\operatorname{Var}(c)$ exceeds the independent value: that comparison would need a separate theorem
relating the average within-$Q$ order-stat variance to the uniform fan. What the identity does isolate
is the between-$Q$ term, one signature of positive dependence (it also shifts the within-$Q$ law from
uniform to $Q$): a per-dataset latent shifts the whole
calibration sample, and the realized coverage rides with it. Empirically the net effect is a wider
fan, e.g.\ a one-factor Gaussian copula at correlation $0.2$ gives $\operatorname{Var}(c)\approx 3.3$
times the independent value (and the comonotone limit $U_1=\dots=U_n$ gives $\tfrac1{12}$, far above
it), but that widening is observed, not proved by Theorem~\ref{thm:pos} alone.

\section{Negative dependence narrows the fan}

The negative regime is the one a genuine prior cannot reach. By \citet{kerns2006} a finite
exchangeable sequence with $\rho<0$ has a \emph{signed} directing measure, so the clean conditioning
of Theorem~\ref{thm:pos} is unavailable: there is no $Q$ to condition on, and the between-dataset
term would be a negative quasi-variance. That is precisely the regime where the fan contracts.

\begin{theorem}[Endpoints]\label{thm:ends}
Independent scores give $c\sim\mathrm{Beta}(k,n-k+1)$, the benchmark fan. At the exchangeable
pairwise-correlation floor, take the discrete rank idealization in which the calibration scores are a
uniformly random permutation of the grid $\{1/(n+1),\dots,n/(n+1)\}$ (sampling without replacement,
pairwise correlation $-1/(n-1)$, the smallest attainable for an exchangeable vector): the order
statistic is deterministic, $U_{(k)}=k/(n+1)$ a.s., so $\operatorname{Var}(c)=0$.
\end{theorem}

Independence is the benchmark, not the maximum: positive dependence can widen the fan well past it
(the comonotone limit reaches $\tfrac1{12}$). The claim of this note is the other direction, that
\emph{negative} association does not widen it, with independence the (conjectural) upper endpoint of
the negative-dependence range. Two model cautions on the floor. The grid permutation has
\emph{discrete}-uniform marginals, not the continuous $\mathrm{Unif}(0,1)$ of $U_i=F(S_i)$, so it sits
in the model only if one allows discrete or rank scores; the continuous-marginal analogue is a Latin
hypercube (one point per stratum, randomly permuted), whose fan is small, of order $1/n^2$, but not
exactly zero. And the endpoints agree with Theorem~\ref{thm:coeff} at $R=1$ and $R=0$. The two are
joined by an exact path.

\begin{proposition}[An exact variance interpolation]\label{prop:path}
Let the calibration law be, with probability $1-t$, $n$ independent uniforms, and with probability
$t$, a contest (a uniformly random permutation of the grid $\{1/(n+1),\dots,n/(n+1)\}$). Both
components yield realized coverage with the same mean $k/(n+1)$, the contest with zero variance, so by
the law of total variance the between-component term vanishes and
\begin{equation}
\operatorname{Var}(c)=(1-t)\,\frac{k(n-k+1)}{(n+1)^2(n+2)} .
\end{equation}
\end{proposition}

\begin{proof}
Condition on the component. $\mathbb E[c\mid\text{iid}]=\mathbb E[c\mid\text{contest}]=k/(n+1)$, so
$\operatorname{Var}(\mathbb E[c\mid\cdot])=0$, and $\mathbb E[\operatorname{Var}(c\mid\cdot)]=
(1-t)\,\tfrac{k(n-k+1)}{(n+1)^2(n+2)}$.
\end{proof}

So the fan collapses linearly to zero as $t\to1$. Two caveats. The mixture's pairwise
correlation is $\operatorname{Corr}(U_i,U_j)=-t/(n+1-2t)$, not $-t/(n-1)$: the grid component
contributes $\operatorname{Cov}=-1/(12(n+1))$ but mixing changes the marginal variance to
$(n+1-2t)/(12(n+1))$; the correlation reaches the floor $-1/(n-1)$ only at $t=1$. And the mixture's
one-dimensional marginal is neither continuous uniform nor discrete-grid uniform, so this is an exact
\emph{variance} interpolation, not a fixed-marginal copula path.

So we can prove the contraction exactly at both endpoints and along this whole interpolating family,
and to leading order everywhere (Theorem~\ref{thm:coeff}). For an \emph{arbitrary} negatively
associated law the contraction also holds once it is summed over all levels.

\begin{theorem}[Aggregate fan bound]\label{thm:agg}
For negatively associated scores with uniform marginals,
\begin{equation}
\sum_{k=1}^n\operatorname{Var}(U_{(k)})\;\le\;\sum_{k=1}^n\frac{k(n-k+1)}{(n+1)^2(n+2)},
\end{equation}
the independent value, with equality holding under independence. Summed over all target levels,
the fan is never widened by negative association.
\end{theorem}

\begin{proof}
The order statistics are a permutation of the sample, so $\sum_k U_{(k)}^2=\sum_i U_i^2$ and, under
uniform marginals, $\sum_k\mathbb E[U_{(k)}^2]=n\,\mathbb E[U_1^2]=n/3$ for every such law; hence
$\sum_k\operatorname{Var}(U_{(k)})=\tfrac n3-\sum_k(\mathbb E\,U_{(k)})^2$. For the top $j$ statistics,
$\sum_{k>n-j}U_{(k)}=\int_0^1\min\!\big(j,n-N(u)\big)\,\mathrm du$. The function $\phi(m)=\max(-j,m-n)$
is convex and \emph{nondecreasing}, which is exactly the class for which \citet{shao2000} compares a
negatively associated sum to its independent counterpart, $\mathbb E\phi(N(u))\le\mathbb E\phi(\mathrm{Bin}(n,u))$.
(Because the means agree, $\mathbb E N(u)=nu=\mathbb E\,\mathrm{Bin}$, one may add any affine term, so
the comparison in fact extends to all convex $\phi$, the full convex order; the monotone-convex form
already suffices here.) Since $\min(j,n-m)=-\phi(m)$, this gives
$\mathbb E\min(j,n-N(u))\ge\mathbb E\min(j,n-\mathrm{Bin}(n,u))$; integrating,
\begin{equation}\label{eq:topj}
\sum_{k>n-j}\mathbb E\,U_{(k)}\;\ge\;\sum_{k>n-j}\tfrac{k}{n+1}\qquad\text{for every }j=1,\dots,n,
\end{equation}
the right side being the independent profile's top-$j$ sum. Both profiles are increasing and have the
same total, $\sum_k\mathbb E\,U_{(k)}=\sum_i\mathbb E\,U_i=n/2=\sum_k\tfrac k{n+1}$ (the case $j=n$ of
\eqref{eq:topj}, with equality). For two increasing vectors with equal sum, ``every top-$j$ partial
sum of one dominates that of the other'' is exactly the definition of majorization; so
$(\mathbb E\,U_{(1)},\dots,\mathbb E\,U_{(n)})$ majorizes $(\tfrac1{n+1},\dots,\tfrac n{n+1})$. The map
$x\mapsto\sum_k x_k^2$ is Schur-convex, hence larger on the majorizing vector,
$\sum_k(\mathbb E\,U_{(k)})^2\ge\sum_k(\tfrac k{n+1})^2$, and the displayed identity for
$\sum_k\operatorname{Var}(U_{(k)})$ finishes it.
\end{proof}

The proof uses only identical uniform marginals and the convex-order domination
$N(u)\le_{\mathrm{cx}}\mathrm{Bin}(n,u)$, not exchangeability as such. Equality holds under
independence; that equality \emph{forces} independence is not established here and would need a
separate tightness lemma.

\section{The single level: false for general negative association}

Everything above is proved for general negative association. The single conformal level, the
operationally relevant one, is not: the per-level bound is
\emph{false} for general NA, at every $n\ge2$. The obstruction the aggregate proof sidesteps is
genuine. The step $\mathbf 1\{N(u)<k\}$ behind a single $\operatorname{Var}(U_{(k)})$ is neither
concave nor convex, so the convex order does not sign it, whereas the smooth aggregate functional
$\min(j,n-N)$ is concave and it does; summation over $k$ is exactly what buys the sign, and a single
level is left unprotected.

\begin{proposition}[Counterexample for every $n$]\label{prop:counter}
For every $n\ge2$ and every small $a>0$ there is an absolutely continuous negatively associated law
with uniform marginals whose top order statistic is strictly over-dispersed,
\[
\operatorname{Var}(U_{(n)})-\operatorname{Var}^{\mathrm{iid}}(U_{(n)})=c_n\,a^2+O(a^3),\qquad c_n>0 ,
\]
with $c_2=\tfrac1{18}$, $c_3=\tfrac3{40}$, $c_4=\tfrac2{25}$, and $c_n>0$ in general. By the negation
symmetry $U\mapsto1-U$ the mirror law over-disperses the minimum $U_{(1)}$ by the same amount.
\end{proposition}

\begin{proof}
Fix $a\in(0,1/n]$, a top bin $T=[1-a,1]$ and bottom bin $B=[0,1-a]$. Draw a category
$Y\in\{0,1,\dots,n\}$ with $\mathbb P(Y=0)=1-na$ and $\mathbb P(Y=i)=a$; set $T_i=\mathbf 1\{Y=i\}$, so
at most one coordinate is ever large; and with independent $V_i\sim\mathrm{Unif}(0,1)$ put
$U_i=(1-a)V_i$ when $T_i=0$ and $U_i=(1-a)+aV_i$ when $T_i=1$. Each marginal is uniform. The law is
negatively associated by closure: $(T_1,\dots,T_n)$ are the cell indicators of a single multinomial
trial on $n+1$ cells, hence NA \citep{joagdev1983}; the $V_i$ are independent and independent of $T$,
so $(T,V)$ is NA; and each $U_i=\varphi(T_i,V_i)$ is nondecreasing in the disjoint pair $(T_i,V_i)$, so
$(U_i)$ is NA. Its diagonal is explicit: $\delta(u)=\mathbb P(U_{(n)}\le u)$ equals
$(1-na)\,(u/(1-a))^n$ on $B$ and $1-n(1-u)$ on $T$ (the union bound is \emph{saturated} there, since at
most one coordinate exceeds $1-a$). With $\mathbb E[U_{(n)}^m]=1-m\int_0^1u^{m-1}\delta(u)\,du$ this
yields the closed form: $\operatorname{Var}(U_{(n)})-\operatorname{Var}^{\mathrm{iid}}$ equals
$a^2(1-6a-2a^2)/18$ at $n=2$ and $a^2(6-64a-45a^2)/80$ at $n=3$, positive for small $a$, and
$c_n\,a^2+O(a^3)$ with $c_n>0$ in general.
\end{proof}

The excess is small but exact: at $n=3$, $a=\tfrac1{20}$ the ratio to the independent fan is $1.0022$;
the violation window is $a\in(0,0.088)$ at $n=3$ and narrows with $n$, but $c_n>0$ for every $n$. Nor is
the extreme level the worst case: a direct search over NA laws over-disperses the \emph{median} order
statistic by more --- about $5\%$ at $n=3$, by three-way repulsion that pins the extremes while letting
the middle coordinate wander --- so the failure is not a boundary artifact of $k=n$. This
is exactly the centroid-violating diagonal the linear program of \eqref{eq:centroid} produces, and a
genuine NA law realises it at every $n$, not merely at $n=2$. The
mechanism is that uniform marginals pin only the \emph{sum} of the order-statistic means (to $n/2$),
never an individual $\mathbb E[U_{(k)}]$, so a mean-decreasing spread is available; the counterexample
lowers $\mathbb E[U_{(n)}]$ while widening it, and the aggregate bound survives because the other levels
over-contract. The excess is necessarily small: negative association \emph{under}-disperses the count
$N(u)$ (Lemma~\ref{lem:count}), whereas a large order-statistic variance wants it over-dispersed, and
the two pressures oppose --- a law whose count is over-dispersed enough to inflate the fan substantially
is positively associated, not NA.

Degeneracy is not the escape either. The failure does not need the counterexample's vanishing
density (it is zero wherever two coordinates are large) or its boundary position in the NA cone (its
binding covariances are exactly zero). It does not: a constrained search produces a law with a
\emph{strictly positive} density (bounded below by $\tfrac15$ of uniform), exactly uniform marginals, and
lying \emph{strictly inside} the NA cone --- every non-trivial disjoint-monotone covariance is
$\le -9.8\times10^{-5}<0$ --- whose smeared continuous form has
$\operatorname{Var}(U_{(3)})=0.037857>\tfrac3{80}$, a ratio $1.0095$. Both facts are certified in exact
rational arithmetic: NA by the complete up-set decision procedure (NA $\iff$
$\operatorname{Cov}(\mathbf 1_S(U_A),\mathbf 1_W(U_B))\le0$ over all disjoint $A,B$ and all up-sets
$S,W$), and the variance by exact integration of the piecewise-uniform smearing (the same routine returns
$\tfrac3{80}$ on the independent law). So full support does not rescue the bound; a positivity floor only
shrinks the achievable excess.

The per-level contraction is therefore simply not a property of negative association --- not for general
NA, not at the extremes, not with a strictly positive density. It \emph{is} a property of the structured
families a calibration design actually produces, and there it is exact or provable.

Three structured families make the point, two of them by outright proof. The
Farlie--Gumbel--Morgenstern copula is the cleanest.

\begin{proposition}[FGM copula]\label{prop:fgm}
Let $(U_1,\dots,U_n)$ have the FGM density $1+\theta\sum_{i<j}(1-2u_i)(1-2u_j)$, valid and negatively
associated for $\theta\in[-1/\binom n2,0]$. Then for every $k$, $\operatorname{Var}_\theta(U_{(k)})$ is
a concave quadratic in $\theta$ with value the independent benchmark $V_0=\eqref{eq:betavar}$ at
$\theta=0$ and nonnegative slope there, hence $\operatorname{Var}_\theta(U_{(k)})\le V_0$, increasingly,
throughout the negative range.
\end{proposition}

\begin{proof}
The density is affine in $\theta$, so $\mathbb E_\theta[U_{(k)}]$ and $\mathbb E_\theta[U_{(k)}^2]$ are
affine in $\theta$ and $\operatorname{Var}_\theta(U_{(k)})=V_0+\beta_k\theta-a_k^2\theta^2$ with
$a_k=\tfrac{d}{d\theta}\mathbb E_\theta[U_{(k)}]$: the quadratic coefficient is minus a square, so the
parabola is concave. Its slope at $0$ is
$\beta_k=\operatorname{Cov}_0\!\big((U_{(k)}-\tfrac k{n+1})^2,\sum_{i<j}(1-2U_i)(1-2U_j)\big)$; by
exchangeability and conditioning on $U_{(k)}=y$ this reduces to
$\beta_k=\tfrac12\,\mathbb E[(Y-\mu)^2N(Y)]$ with $Y\sim\mathrm{Beta}(k,n{-}k{+}1)$, $\mu=k/(n+1)$ and
$N(y)=(n-1)(n+2)y^2-2(kn-1)y+k(k-1)$, which evaluates to
$\beta_k=k(n{+}1{-}k)\,Q/[(n+1)^2(n+2)(n+3)(n+4)]$ with $Q=(n-11)k(n{+}1{-}k)+(3n+2)(n+1)$. Writing
$t=k(n{+}1{-}k)\in[n,(n+1)^2/4]$: for $n\ge11$, $Q\ge(3n+2)(n+1)>0$; for $n\le10$,
$Q\ge(n-11)\tfrac{(n+1)^2}4+(3n+2)(n+1)=(n-1)(n+1)(n+3)/4\ge0$. So $\beta_k\ge0$, and with concavity the
variance is nondecreasing and $\le V_0$ on $[-1/\binom n2,0]$.
\end{proof}

Sampling designs are the practically central case --- they are the literal \emph{design of the
calibration sample} --- and there the contraction is exact, with a closed form, at every level.

\begin{proposition}[Sampling without replacement]\label{prop:wor}
Draw $n\le M$ scores without replacement from the regular grid population $\{r/(M+1):r=1,\dots,M\}$ (the
multivariate hypergeometric, NA by \citealp{joagdev1983}). Then for every $k$,
\[
\mathbb E[U_{(k)}]=\frac{k}{n+1},\qquad
\operatorname{Var}(U_{(k)})=\frac{M-n}{M+1}\cdot\frac{k(n-k+1)}{(n+1)^2(n+2)}\;\le\;\frac{k(n-k+1)}{(n+1)^2(n+2)} .
\]
\end{proposition}

\begin{proof}
The rank $R_{(k)}$ of the $k$-th order statistic is negative-hypergeometric; writing
$R_{(k)}=k+\sum_{i<k}G_i$ with $(G_0,\dots,G_n)$ the exchangeable composition of the $M-n$ unsampled
points into $n+1$ gaps gives $\mathbb E R_{(k)}=k(M+1)/(n+1)$ and $\operatorname{Var}(R_{(k)})=
\tfrac{k(n-k+1)}{(n+1)^2(n+2)}(M+1)(M-n)$; divide by $M+1$ and $(M+1)^2$.
\end{proof}

The finite-population factor $\tfrac{M-n}{M+1}\in[0,1)$ shrinks \emph{every} level by the same amount,
with the coverage mean exactly on target; it runs to $1$ as $M\to\infty$ (WOR $\to$ i.i.d.) and to $0$ at
the census $M=n$ (the deterministic grid, zero fan). Two relatives are equally clean. \emph{Stratified sampling} ---
$L$ equal strata, $m$ i.i.d.-uniform points each, indices permuted (NA by independent union and
permutation, \citealp{joagdev1983}) --- gives $\operatorname{Var}(U_{(k)})=V_0(m,s)/L^2\le V_0(n,k)$ for
$k=(l-1)m+s$, where $V_0(a,b)=\tfrac{b(a-b+1)}{(a+1)^2(a+2)}$; the ratio is $(L-1)$ times a nonnegative
cofactor. \emph{Latin hypercube} sampling (the $m=1$ case, one point per equal-probability stratum) gives
the exact $\operatorname{Var}(U_{(k)})=1/(12n^2)$, an order below the i.i.d.\ $O(1/n)$ at every level, and
conservative on the mean. So the designs a practitioner actually builds all contract, and WOR does so
without shifting the center. One caution: this needs a genuinely fixed, spread population; subsampling a
\emph{random i.i.d.\ pool} returns i.i.d.\ scores and no contraction. Two further natural NA structures
join them numerically: uniform spacings / $\mathrm{Dirichlet}(1,\dots,1)$ shares and Thurstone / Luce
contest shares contract at every level (ratios $0.27$--$0.74$), conservatively at the operating quantile.

And in the canonical continuous family the per-level statement is exact, and more, monotone in the
correlation.

\begin{proposition}[Exact per-level law for equicorrelated normals]\label{prop:gauss}
Let $(Z_1,\dots,Z_n)$ be equicorrelated standard normal with correlation
$\rho\in[-\tfrac1{n-1},1]$, and let $v_k=\operatorname{Var}(\xi_{(k)})$ be the variance of the $k$-th
order statistic of $n$ i.i.d.\ standard normals. Then, for every $k$,
\begin{equation}
\operatorname{Var}(Z_{(k)})=v_k+\rho\,(1-v_k).
\end{equation}
As $v_k<1$ for $n\ge 2$, this is strictly increasing in $\rho$: negative correlation strictly shrinks
the variance of every order statistic, positive correlation inflates it.
\end{proposition}

\begin{proof}
Write $Z_i=\bar Z+(Z_i-\bar Z)$. For equicorrelated normals $\bar Z\sim N(0,s)$,
$s=(1+(n-1)\rho)/n$, is independent of the centred vector, and
$\operatorname{Cov}(Z_i-\bar Z,Z_j-\bar Z)=(1-\rho)(\delta_{ij}-1/n)$, exactly $1-\rho$ times the
covariance of $(\xi_i-\bar\xi)$ for i.i.d.\ standard normals. Hence
$Z_{(k)}\stackrel{d}{=}\sqrt{s}\,\eta+\sqrt{1-\rho}\,(\xi-\bar\xi)_{(k)}$ with $\eta\sim N(0,1)$
independent of the second term, so
$\operatorname{Var}(Z_{(k)})=s+(1-\rho)\operatorname{Var}\big((\xi-\bar\xi)_{(k)}\big)$. Both terms
are affine in $\rho$; matching $v_k$ at $\rho=0$ forces
$\operatorname{Var}((\xi-\bar\xi)_{(k)})=v_k-1/n$, and substituting gives $v_k+\rho(1-v_k)$.
\end{proof}

\begin{remark}
This is the per-level contraction, exact and monotone, on the Gaussian scale, unifying the two sides:
the single line $v_k+\rho(1-v_k)$ runs from inflation at $\rho>0$ through the independent value at
$\rho=0$ to the floor at $\rho=-1/(n-1)$. Under the Gaussian copula $U_i=\Phi(Z_i)$, which restores
uniform marginals, $\operatorname{Var}(U_{(k)})$ is numerically monotone in $\rho$ and below the
independent value throughout $\rho<0$ (\texttt{check\_fan.py}). The transfer to the uniform scale is
a dispersive-order statement, and it is half rigorous. Decompose along the common mean,
$Z_{(k)}=\sqrt{1-\rho}\,D+\bar Z$ with $D=(\xi-\bar\xi)_{(k)}$ and $\bar Z\sim N(0,s)$ independent.
$D$ is log-concave, being the order statistic of a vector with log-concave (Gaussian) joint density
\citep{prekopa1973}; and convolution with a log-concave density increases the dispersive order
\citep{lewis1981}, so adding the Gaussian shift $\bar Z$ disperses, monotonically in $s$. The residual
is that this outweighs the simultaneous $\sqrt{1-\rho}$ shrinkage of $D$ as $\rho$ grows, a one-parameter
density-quantile monotonicity under the variance budget $\tfrac{n-1}n(1-\rho)+s=1$, which no
off-the-shelf convolution theorem delivers but which we now settle.

Tweedie's formula sharpens that residual to a single posterior-variance bound. Reading
$Y:=Z_{(k)}=aD+b\eta$ ($a=\sqrt{1-\rho}$, $b=\sqrt s$) as Gaussian signal extraction of the
log-concave signal $D$, the slope of the posterior mean equals the scaled posterior variance,
$\frac{d}{dq}\mathbb E[D\mid Y=q]=a\,\operatorname{Var}(D\mid Y=q)/b^2$
\citep{efron2011,hansen2026}. The dispersive monotonicity is equivalent to this slope being at most
$(n-1)a/n$, so, using $\tfrac{n-1}n a^2+b^2=1$, the entire transfer reduces to
\begin{equation}
\operatorname{Var}(D\mid Y=q)\;\le\;\frac{n-1}{n}\,b^2\qquad\text{for all }q,
\end{equation}
a posterior-variance bound for log-concave Gaussian denoising, which the next lemma settles. The
score-driven reading of the same identity \citep{hansen2026} is the econometric form of modelling the
conditional spread, which is exactly what closes the information gap a conformal wrapper leaves
\citep{cotton_marginally}.
\end{remark}

\begin{lemma}[Strong log-concavity of $D$]\label{lem:slc}
$D=(\xi-\bar\xi)_{(k)}$ is $\tfrac{n}{n-1}$-strongly log-concave, $(\log g_D)''\le -n/(n-1)$.
\end{lemma}

\begin{proof}
The centred vector $w=\xi-\bar\xi$ has covariance $I-\tfrac1n J$, an orthogonal projection, so on the
hyperplane $\{\sum_i w_i=0\}$ it is a \emph{standard} Gaussian: its $-\log$ density is $\tfrac12\|w\|^2$,
with Hessian the identity. The law of the ordered vector is this density restricted to the convex cone
$\{w_1\le\dots\le w_n\}$. The cone enters only through the indicator of a convex set, a log-concave
factor; writing the restricted density as $\exp(-\tfrac12 w^\top(I-\tfrac1nJ)^{+}w-\iota_C(w))$ with
$\iota_C$ the convex indicator of the cone, the Hessian of the potential is
$(I-\tfrac1nJ)^{+}+\nabla^2\iota_C\succeq(I-\tfrac1nJ)^{+}$. So the strong-log-concavity constant is
set by the Gaussian quadratic alone; the cone can only raise it, never lower it, and in particular
does not interfere through its boundary. Strong log-concavity passes to marginals
\citep{prekopa1973,brascamplieb1976,saumard2014}, the marginal of a coordinate being strongly
log-concave to the reciprocal of that coordinate's reference variance; for $w_{(k)}$ this is
$\tfrac{n-1}n$ (the diagonal of $I-\tfrac1nJ$), giving constant at least $\tfrac n{n-1}$ (the
parameter bookkeeping is in the Appendix).
\end{proof}

\begin{proposition}[Gaussian-scale dispersive monotonicity]\label{prop:dispmon}
For equicorrelated standard normals with $\rho\in(-\tfrac1{n-1},1)$, $Z_{(k)}$ is increasing in the
dispersive order as $\rho$ increases.
\end{proposition}

\begin{proof}
For a smooth one-parameter family the quantile evolves by $\partial_\rho q_\rho(u)=
\mathbb E[\dot Y\mid Y=q_\rho(u)]$ with $\dot Y=\partial_\rho Z_{(k)}$, and dispersive monotonicity is
exactly $\partial_\rho q_\rho(u)$ increasing in $u$, i.e.\ increasing in $q$. Writing
$Y=aD+b\eta$ and using $a\,\mathbb E[D\mid Y]+b\,\mathbb E[\eta\mid Y]=Y$ together with the budget
$\tfrac{n-1}na^2+b^2=1$, this reduces to $m'(q)\le (n-1)a/n$, where $m(q)=\mathbb E[D\mid Y=q]$ (the
computation of $M'(q)$ is in the Appendix).
Tweedie's identity for Gaussian signal extraction gives $m'(q)=a\operatorname{Var}(D\mid Y=q)/b^2$.
The posterior of $D$ given $Y$ has $-\log$ density with Hessian
$(-\log g_D)''+a^2/b^2\ge \tfrac n{n-1}+\tfrac{a^2}{b^2}=\tfrac{n}{(n-1)b^2}$, the last equality by the
budget and Lemma~\ref{lem:slc}; the Brascamp--Lieb variance inequality \citep{brascamplieb1976} then
gives $\operatorname{Var}(D\mid Y=q)\le (n-1)b^2/n$, hence $m'(q)\le (n-1)a/n$. Equivalently the
quantile density $q_\rho'(u)=1/f_\rho(q_\rho(u))$ increases in $\rho$ at each level (the density at each
quantile decreases), which is the dispersive monotonicity. At the open-interval ends the statement is
read in the limit: $b\to0$ at $\rho=-1/(n-1)$ (the degenerate floor, $Z_{(k)}\to\sqrt{n/(n-1)}\,D$) and
$a\to0$ at $\rho=1$ ($Z_{(k)}\to\eta$).
\end{proof}

This settles the Gaussian-scale per-level statement in the strong (dispersive, not merely variance)
sense. The passage to the uniform scale follows by differentiating the variance directly in the
correlation matrix. For $Z\sim N(0,R)$ with $R_{ij}\le0$ (NA), along the path $R(t)=(1-t)R+tI$ to independence,
Price's theorem gives $\tfrac{d}{dt}\operatorname{Var}(h(Z_{(k)}))=\sum_{i<j}(-R_{ij})D_{ij}$ for any
$C^1$ map $h$, with the closed form (one Gaussian integration by parts)
$D_{ij}=2\,\mathbb E[(h(Z_{(k)})-\mu_h)\,h'(Z_{(k)})\,\mathbf 1\{i=\mathrm{rank}_k\}\,(R^{-1}Z)_j]$. Taking
$h=\Phi$ makes the endpoint at $R=I$ \emph{exactly} the $\mathrm{Beta}(k,n-k+1)$ fan, so no order need
survive $\Phi$; the whole statement, on both scales and for arbitrary nonpositive off-diagonal $R$,
reduces to the single sign $D_{ij}\ge0$. That sign is proved for $n=2$ and on the equicorrelated slice
(where it recovers $v_k+\rho(1-v_k)$), and certified across the negative cone to $n=8$; the remaining
crux is an adjacent-tie-surface domination.

The Gaussian case fixes the copula and varies $k$; the opposite slice, the top level $k=n$ for
\emph{any} negatively associated sample, is also exactly tractable and is where the conjecture becomes
one-dimensional. There $U_{(n)}\le u$ iff every score lies below $u$, so the law of $U_{(n)}$ is the
copula diagonal $\delta(u)=C(u,\dots,u)=\mathbb P(N(u)=n)$, and
$\operatorname{Var}(U_{(n)})=\int_0^1\!\!\int_0^1[\delta(u\wedge v)-\delta(u)\delta(v)]\,du\,dv$.
Negative association enters only as the negative-orthant bound $\delta(u)\le u^n$
\citep{joagdev1983} --- equivalently the convex-order image of the discretely convex indicator
$\mathbf 1\{N=n\}$. Writing the deficit $g(u)=u^n-\delta(u)\ge 0$ and $\mu^\ast=n/(n+1)=\mathbb
E[U_{(n)}^{\mathrm{iid}}]$, the variance has the exact closed form
\begin{equation}
\operatorname{Var}(U_{(n)})=\operatorname{Var}^{\mathrm{iid}}(U_{(n)})
+2\int_0^1(u-\mu^\ast)\,g(u)\,du-\Big(\int_0^1 g\Big)^{\!2},
\end{equation}
so the per-level bound at $k=n$ is \emph{equivalent} to a centroid inequality on the deficit,
\begin{equation}\label{eq:centroid}
\frac{\int_0^1 u\,g(u)\,du}{\int_0^1 g(u)\,du}\;\le\;\mu^\ast+\tfrac12\!\int_0^1 g\,:
\end{equation}
the mass by which the diagonal falls short of independence must sit, on average, at most a half-deficit
above the independent-maximum mean. This is \emph{not} a consequence of the diagonal alone. A linear
program over the functions meeting every condition a copula diagonal must satisfy (monotone, $0\to1$,
$\delta\le u$, $n$-Lipschitz, Fr\'echet lower bound; \citealp{nelsen2006}) together with the negative-orthant
bound $\delta\le u^n$ returns a diagonal that violates \eqref{eq:centroid} (\texttt{check\_fan.py}).
The single level is therefore irreducibly a property of the off-diagonal negative-orthant probabilities
across thresholds, not of the diagonal that determines $U_{(n)}$'s law: exactly the content the
aggregate convex-order bound uses in bulk and that a one-level statement cannot read off. It is the same
obstruction that blocks localizing Theorem~\ref{thm:agg} to a single $k$, now in closed form. The
LP-violating diagonal is realised by the genuine NA law of Proposition~\ref{prop:counter} at every
$n$ --- its top-bin diagonal $1-n(1-u)$ saturates the union bound, front-loading the deficit onto the
bottom bin with centroid above $n/(n+1)$ --- so \eqref{eq:centroid} genuinely fails, and the bound is
recovered only on the structured families of Propositions~\ref{prop:gauss}--\ref{prop:fgm} and
without-replacement sampling.

The same convex-order contraction $N(u)\le_{\mathrm{cx}}\mathrm{Bin}(n,u)$ that drives the aggregate
bound is available per level, but turning it into the single-$k$ statement needs the tail
$\mathbb P(N(u)\ge k)$ to cross the binomial tail once in $u$, which the total positivity of the
binomial tail in $(k,u)$ \citep{karlin1968} makes plausible, followed by a
single-crossing-implies-less-spread step. We claim only the variance statement, and deliberately not
the stronger dispersive ordering: an order statistic of a dependent sample need not be dispersively
smaller than the independent one even at fixed marginals \citep{panja2022}, so the result is specific
to the variance and is not a corollary of a dispersive theorem. The natural sufficient condition on
the dependence is negative supermodular dependence \citep{christofides2004}, which all the structures
below satisfy. The table is illustrative, each row a single Monte Carlo estimate ($6\times10^4$ draws,
so the variance entries carry a relative Monte Carlo error of a few percent), fully reproduced by
\texttt{check\_fan.py} ($n=20$, $\alpha=0.1$, independent fan variance $3.92\times10^{-3}$); it is
evidence for the conjecture, not a proof.

\begin{center}
\begin{tabular}{lcccc}
\hline
calibration dependence & $\mathbb E[c]$ & $\operatorname{Var}(c)$ & ratio to fan & sign \\
\hline
independent ($\rho=0$)                 & $0.904$ & $3.92\times10^{-3}$ & $1.00$ & --- \\
exchangeable, $\rho=-1/(n-1)$          & $0.923$ & $1.16\times10^{-3}$ & $0.29$ & neg.\ assoc. \\
MA(1), lag-1 corr $-\tfrac12$          & $0.912$ & $2.71\times10^{-3}$ & $0.69$ & neg.\ assoc. \\
random negatively-paired blocks        & $0.909$ & $3.14\times10^{-3}$ & $0.80$ & neg.\ assoc. \\
contest / without replacement          & $0.905$ & $0$                  & $0.00$ & floor \\
one-factor, $\rho=+0.2$                & $0.863$ & $1.33\times10^{-2}$ & $3.41$ & positive \\
\hline
\end{tabular}
\end{center}

Every negatively associated structure, exchangeable or not, comes in under the independent fan; the
positive control widens it. The mean of $c$ also responds to dependence, drifting up under negative
and down under positive (to $0.863$ in the one-factor row), confirming it is not pinned; the variance
is what we track here. A wider search agrees:
over exchangeable, moving-average, block, and randomly-generated non-exchangeable negatively-correlated
Gaussian copulas at $n\le 8$ and every $k$, the ratio to the independent fan never exceeds one
(it ranges from about $0.30$ at the floor to $0.80$ for weak dependence), with no counterexample. In all
of these the order-statistic CDFs $\mathbb P(N(u)\ge k)$ cross the independent ones exactly once, the
single crossing the remark below uses.

\begin{remark}[A closed form, and what an exact proof must dominate]
The variance is an explicit functional of $g(u)=\mathbb P(N(u)\ge k)$,
\begin{equation}
\operatorname{Var}(c)=2\int_0^1(1-u)\,g(u)\,\mathrm du-\Big(\int_0^1 g(u)\,\mathrm du\Big)^{\!2},
\end{equation}
so writing $\delta=g-g^{\mathrm{iid}}$, $A=\int_0^1 g^{\mathrm{iid}}=1-\mathbb E[c^{\mathrm{iid}}]$,
and $\Delta=\int_0^1\delta$,
\begin{equation}
\operatorname{Var}(c)-\operatorname{Var}(c^{\mathrm{iid}})
=2\int_0^1\big[(1-u)-A\big]\,\delta(u)\,\mathrm du-\Delta^2 .
\end{equation}
The $-\Delta^2$ term is non-positive. The weight $(1-u)-A$ changes sign at
$u=\mathbb E[c^{\mathrm{iid}}]=k/(n+1)$. Empirically (across every structure in the table, and in the
wider search below) $\delta$ changes sign exactly once, near $u\approx k/n$; this is an
\emph{observed} single-crossing, not a consequence of Lemma~\ref{lem:count}, since count-variance
contraction alone does not force one-crossing of the tail curve $g$. \emph{Granting} that single
crossing, the two sign changes agree up to the conformal offset between $k/n$ and $k/(n+1)$, a band of
width $O(1/n)$ on which the integrand is positive and elsewhere negative, so the identity gives the
contraction up to that $O(1/n)$ boundary term, a second route to the leading order of
Theorem~\ref{thm:coeff}, and isolates what a finite-sample proof must control: the boundary remainder
against the slack $\Delta^2$.
\end{remark}

\begin{remark}[Why the closed form does not close it, and the lemma that would]
The closed form is, by itself, circular. Substituting the moment identities
$\int u\,\delta=\tfrac12(\mathbb E[c_{\mathrm{iid}}^2]-\mathbb E[c^2])$ and
$\int\delta=\mathbb E[c_{\mathrm{iid}}]-\mathbb E[c]$ into any elementary rearrangement of the
difference, single-crossing splits, Abel summation, the exact first-order perturbation, collapses it
back to $\operatorname{Var}(c)-\operatorname{Var}(c_{\mathrm{iid}})$ with the auxiliary crossing
point cancelling: these structural facts carry no magnitude. Nor is the per-level convex order of
Lemma-\ref{lem:count} enough on its own. Because $\operatorname{Var}(c)$ depends only on the single
curve $g(u)=\mathbb P(N(u)\ge k)$, the pointwise convex order $N(u)\le_{\mathrm{cx}}\mathrm{Bin}(n,u)$
does not suffice, since a curve dominated pointwise can have larger variance (place an atom
at the boundary), and such curves are excluded only because they are not realizable by a genuine
sample. The content is therefore an extremal statement about the realizable curves, not their
marginals:
\begin{quote}
\emph{False in general.} Among all $g(u)=\mathbb P(N(u)\ge k)$ realizable by a negatively associated
uniform sample, the independent (binomial) curve does \emph{not} maximize $\operatorname{Var}(c)$:
Proposition~\ref{prop:counter} realizes a larger one at every $n$. The extremal curve is maximized by
the independent one only within structured subclasses.
\end{quote}
What a total-positivity or coupling argument delivers is thus not the universal bound but its
restriction to a copula family; Proposition~\ref{prop:fgm} is the clean instance, and
Proposition~\ref{prop:path} and the contest endpoint are exact cases.
\end{remark}

\begin{remark}[The signed measure as a negative between-term]
Theorem~\ref{thm:pos} reads the fan as within-dataset variance plus a non-negative between-dataset
variance. In the signed-de-Finetti picture of the companion note, the negative corner is exactly
where that between-dataset object turns negative, a Feynman--Wigner quasi-variance that subtracts from
the within-dataset fan rather than adding to it, reaching the full cancellation $R=0$ at the contest
floor. This is heuristic, since a signed measure has no genuine conditional variance, but it is the
right cartoon: the same sign that decides conditional adaptivity in \citet{cotton_signed} decides
whether the fan is inflated or cancelled here.
\end{remark}

\begin{remark}[The positive side, and an asymmetry]
Positive dependence --- the common case, since serial dependence in a time series, clustered or grouped
calibration data, and shared latent factors all produce it --- runs the mirror direction: the aggregate
fan never \emph{narrows} (the majorization of Theorem~\ref{thm:agg} reverses, $N(u)\ge_{\mathrm{cx}}
\mathrm{Bin}$), and every structured family widens \emph{every} level. For equicorrelated normals
Proposition~\ref{prop:gauss} is exact and monotone: $\rho>0$ inflates each $\operatorname{Var}(Z_{(k)})$;
a one-factor Gaussian copula at $\rho=0.2$ gives $\approx 3.3\times$ the independent fan, an AR$(1)$
calibration segment at lag-$1$ correlation $0.6$ gives $\approx 2.5\times$, and the comonotone limit
reaches $\tfrac1{12}$. Whether a single level can buck this --- narrowing below the benchmark the way
the negative side widens above it --- is open. The obstruction is symmetric --- association gives
$\operatorname{Cov}(\mathbf 1\{N(u)\ge k\},\mathbf 1\{N(v)\ge k\})\ge0$ but not the comparison to the
independent value --- yet an extensive exact search (the up-set positive-association decision procedure,
smeared variance validated against the benchmark) finds \emph{no} narrowing law, and a descent over the
positive-association region floors above the benchmark at every level ($1.09$ at the median).
\begin{conjecture}[No positive buck]\label{conj:posmono}
If the calibration scores are positively associated with uniform marginals, then
$\operatorname{Var}(U_{(k)})\ge$ the independent value \eqref{eq:betavar} for every $k$.
\end{conjecture}
\noindent If it holds, the fan is treacherous exactly where it would help --- under negative dependence a
single level can widen against the aggregate trend --- and dependable exactly where it hurts, so the
long-run-variance inflation of the next section is a \emph{safe} pessimistic correction.
\end{remark}

\section{A constructive corollary}

If negative dependence narrows the fan, one can engineer it, and there is a clean continuous-marginal
construction that stays inside the $U_i\sim\mathrm{Unif}(0,1)$ model (unlike the discrete grid of
Theorem~\ref{thm:ends}). Take a one-per-stratum design: with $\pi$ a uniform random permutation of
$\{1,\dots,n\}$,
\begin{equation}
U_i=\frac{\pi(i)-1+V_i}{n}\quad(\text{Latin hypercube, }V_i\stackrel{\text{iid}}{\sim}\mathrm{Unif}(0,1)),
\qquad\text{or}\qquad U_i=\frac{\pi(i)-1+T}{n}\quad(\text{common jitter, }T\sim\mathrm{Unif}(0,1)).
\end{equation}
Each $U_i$ is exactly $\mathrm{Unif}(0,1)$ and the sample is exchangeable, with one point per stratum,
so the $k$-th order statistic lies in the $k$-th stratum and
$\mathbb E[U_{(k)}]=(k-\tfrac12)/n$, $\operatorname{Var}(U_{(k)})=1/(12n^2)$, a fan of order $1/n^2$
rather than $1/n$. (Note these continuous designs give mean $(k-\tfrac12)/n$, not the grid's exact
$k/(n+1)$.) More generally the classical variance-reduction designs apply: stratification never
increases variance \citep{cochran1977}, antithetic pairing helps for the monotone empirical CDF, and
determinantal or repulsive designs give negative association with faster-than-independent
concentration \citep{hough2006,bardenet2016}.

\begin{remark}[Admissible designs]
The marginal guarantee needs each calibration score exchangeable with the test score, not the
calibration scores independent of one another, so negative association is in principle compatible. But
the routes differ in what they cost. Score-dependent curation of the calibration set, the obvious way
to balance it, breaks exchangeability with the test point and voids the guarantee. A pooling that uses
all $n+1$ scores together (pooled ranks or de-meaning) preserves rank validity, but it is
transductive, full-conformal in flavour, not split conformal with a threshold fixed before the test
point is seen, and once the test score enters the pooled construction the realized-coverage variable
is no longer the fixed-threshold $c=U_{(k)}$ studied here. The safe message is the design one: a
pre-specified, exchangeability-preserving randomized design can lower the run-to-run variance of the
threshold without disturbing validity. What makes such a design admissible is not that it is merely
randomized and pre-specified, but that it preserves the exchangeability the validity argument needs;
covariate stratification qualifies only when it is external to the responses and applied symmetrically
across the $n+1$ points. Exact zero fan is the discrete rank/grid idealization; continuous-marginal
designs (Latin hypercube) give a small but nonzero fan, of order $1/n^2$. The tension is intrinsic: a
generic negatively associated design lowers the dispersion only by also nudging $\mathbb E[c]$ off
$1-\alpha$ (the mean is not pinned), so coverage that is both reproducible \emph{and} on-target demands
a \emph{mean-preserving} symmetric design. Distribution-free validity and minimal coverage dispersion
cannot both be driven to their limit at once, save in the degenerate zero-fan grid.
A deeper caveat: the dispersion is a \emph{within-model} quantity, computed with the marginal $F$ and
the copula $C$ held fixed. Shrinking it by design lowers sampling variability but does nothing for model risk (whether $F$ is the right target at all) and removes the run-to-run scatter that is
the cheapest diagnostic of misspecification. A vanishing fan against an assumed $F$ certifies
reproducibility, not correctness; the zero-fan grid reads $1-\alpha$ every draw by construction, not
because the world obliges.
\end{remark}

\begin{remark}[The copula viewpoint]\label{rmk:copula}
Because the scores are read on the uniform PIT scale, the joint law of $(U_1,\dots,U_n)$ is, by Sklar's
theorem, a copula $C$, with the marginal $F$ factored out entirely. So $\operatorname{Var}(c)$ is a
\emph{copula functional}: the fan is a copula invariant, identical for any score whose calibration
dependence is $C$. The fan coefficient is read off the bivariate copula diagonal,
$\bar c=\tfrac1{n(n-1)}\sum_{i\ne j}\big(C_{ij}(p,p)-p^2\big)$ at the operating quantile $p=1-\alpha$;
for the maximum ($k=n$) $\operatorname{Var}(c)$ is exactly a functional of the diagonal $u\mapsto
C(u,u)$, since $\mathbb P(U_{(n)}\le u)=C(u,u)$. The endpoints are the Fr\'echet bounds: independence
$\Pi$ gives the Beta fan, the comonotone copula $M$ ($U_1=\dots=U_n$) the maximal fan $\tfrac1{12}$;
the negative side has no dual, as the countermonotone lower bound is not a copula for $n>2$, which is
why the floor is only the most negative \emph{exchangeable} copula, $\rho=-1/(n-1)$.
Theorem~\ref{thm:pos} then says that the \emph{extendable} copulas, the exchangeable Archimedean
families (Clayton, Gumbel, Frank), each with a de Finetti frailty $W$ as directing variable, carry a
non-negative between-$Q$ term and widen the fan; the negative regime is the signed-de-Finetti corner
with no such mixture, which is why its contraction needed convex order rather than conditioning.
Finally, tail dependence acts through the copula's tails: a copula dependent in one tail pins the
order statistics there and fattens the opposite end, so lower-tail dependence (Clayton) inflates the
high, conformal quantile more than upper-tail dependence (Gumbel) --- the reverse of the naive guess,
as the companion demonstration shows.
\end{remark}

\section{The time-series fan}\label{sec:tsfan}

When the calibration scores are a stationary segment of a series (the deployment mode priced by \citealp{barberpananjady2026}), the coupled-test caution of the setup applies: the coverage
statement for the next step of the same series is theirs, not ours. What the machinery above
supplies is the dispersion of the \emph{threshold} $U_{(k)}$ across calibration segments, and
with it the second moment of realized coverage against the stationary law. The mean shifts by at most the switch-coefficient tax,
$|\mathbb E[c]-(1-\alpha)|$ bounded through $\min_\tau\{\tau/(n+1)+2\beta(\tau)\}$ by their two-sided
bounds. The variance is governed by the same covariance kernel as Theorem~\ref{thm:coeff}: in the
Bahadur regime,
\[
\operatorname{Var}(c)\;=\;\frac{\sigma^2_{\mathrm{LR}}}{n}+o(1/n),
\qquad
\sigma^2_{\mathrm{LR}}(u)\;=\;\sum_{j\in\mathbb Z}\operatorname{Cov}\bigl(\mathbf 1\{U_0\le
u\},\,\mathbf 1\{U_j\le u\}\bigr),
\]
the long-run variance of the sub-threshold indicator process at the operating quantile ---
equivalently $2\pi$ times its spectral density at zero. For positively serially dependent scores
$\sigma^2_{\mathrm{LR}}\ge\alpha(1-\alpha)$ and the fan widens by the constant factor
$\sigma^2_{\mathrm{LR}}/(\alpha(1-\alpha))$, which does not shrink with $n$. Realized coverage on
a time-series calibration set is thus approximately
$\mathcal N\!\bigl(1-\alpha-\text{tax},\ \sigma^2_{\mathrm{LR}}/n\bigr)$: the mean damage
vanishes at rate $1/n$ while the dispersion carries a permanent inflation factor. Dependence
therefore presents three separate charges: the vanishing mean tax, the fixed fan inflation, and the fixed sharpness rent of the companion note \citep{cotton_twoprices}. Only the first is reduced by more calibration data.

\section{Related work}

The $\mathrm{Beta}(k,n-k+1)$ law of calibration-conditional coverage and its variance are
\citet{vovk2012} and \citet{marques2024}; to our knowledge the papers cited here do not study the
dependence among calibration scores as a control variable for the realized-coverage variance. The
loss-side generalization of the calibration-conditional law is \citet{snellgriffiths2025}, who put
$\mathrm{Dirichlet}(1,\dots,1)$ weights on the ordered calibration losses via the
Aitchison--Dunsmore quantile-spacing result and report a posterior over realized risk: with the
0--1 miscoverage loss their construction collapses to the Beta fan above, their recovery of
conformal risk control is the fan's \emph{mean}, and their highest-posterior-density rule is the
prescription ``report the fan's quantile, not its mean.'' Their guarantee is what they call data-conditional (conditional on the calibration draw, our axis) and, as they note, not input-conditional; it moves along the fan, not across the information gap. Negative association and its consequences are \citet{joagdev1983}; that negatively associated
sums are dominated in convex order by independent ones is \citet{shao2000}, and sampling without
replacement concentrates at least as tightly as with replacement \citep{serfling1974}. The total
positivity behind the single-crossing step is \citet{karlin1968}; that an order statistic of a
dependent sample need not be dispersively smaller than the independent one, which is why we claim
only the variance, is \citet{panja2022}; the closest prior comparison of order statistics across
dependence, \citet{huhu1998}, orders their \emph{distribution functions} between associated and
independent samples rather than the \emph{variance} that is our object; the connection between
supermodular ordering and negative dependence is \citet{christofides2004}; the Bahadur representation
is \citet{bahadur1966}. Determinantal and repulsive designs are \citet{hough2006,bardenet2016} and
stratification is \citet{cochran1977}; \citet{wieczorek2023} uses ``design-based'' conformal
prediction in the survey-sampling sense---randomization-based validity under a complex sampling
design---distinct from the variance-reduction designs here. The finite signed de Finetti representation is
\citet{kerns2006}, and the within-dataset, across-$x$ reading of its sign is the companion note
\citep{cotton_signed}; the information gap is \citep{cotton_marginally}. The de Finetti directing
measure behind the positive side is a prior in all but name; the Bayesian--conformal interface, where
priors meet finite-sample coverage and the trade is validity against efficiency, is surveyed by
\citet{deliu2025}, and the dispersion-versus-guarantee tension above is a variance-side instance of it.
The contribution here is to
make the dependence the control variable for the \emph{variance} of realized coverage: the
leading-order fan coefficient, the exact law-of-total-variance decomposition on the positive side, the
exact aggregate contraction under negative association (Theorem~\ref{thm:agg}), and, for equicorrelated
normals, the exact and dispersively-monotone per-level law. The general single-level negative-association
contraction is left as a conjecture. The copula reading (Remark~\ref{rmk:copula}) rests on Sklar's
theorem \citep{sklar1959,nelsen2006}; copulas enter conformal prediction elsewhere through the
multivariate \emph{output} (e.g.\ \citealp{sunyu2024}), a different axis from the calibration-set
copula here.

\section*{Appendix: the quantile-slope reduction for Proposition~\ref{prop:dispmon}}

With $Z_{(k)}=aD+b\eta$, $a=\sqrt{1-\rho}$, $b=\sqrt{(1+(n-1)\rho)/n}$, the budget
$(n-1)a^2+nb^2=n$ holds and $da/d\rho=-1/(2a)$, $db/d\rho=(n-1)/(2nb)$. Under the natural coupling
(the same $D,\eta$ for all $\rho$), $\dot Y:=\partial_\rho Z_{(k)}=-D/(2a)+(n-1)\eta/(2nb)$, and the
quantile evolves by $\partial_\rho q_\rho(u)=\mathbb E[\dot Y\mid Y=q]=:M(q)$. Writing
$m(q)=\mathbb E[D\mid Y=q]$ and $b\,\mathbb E[\eta\mid Y=q]=q-a\,m(q)$,
\begin{equation}
M(q)=-\frac{m(q)}{2a}+\frac{n-1}{2nb^2}\big(q-a\,m(q)\big).
\end{equation}
Differentiating and using $\tfrac1{2a}+\tfrac{(n-1)a}{2nb^2}=\tfrac{nb^2+(n-1)a^2}{2anb^2}=\tfrac1{2ab^2}$
(the budget),
\begin{equation}
M'(q)=\frac{n-1}{2nb^2}-\frac{m'(q)}{2ab^2},
\end{equation}
so $M'(q)\ge0$, dispersive monotonicity ($\partial_\rho q_\rho$ increasing in $q$), is equivalent to
$m'(q)\le (n-1)a/n$. Tweedie's identity $m'(q)=a\operatorname{Var}(D\mid Y=q)/b^2$ turns this into the
posterior-variance bound $\operatorname{Var}(D\mid Y=q)\le (n-1)b^2/n$, which Lemma~\ref{lem:slc} and
the Brascamp--Lieb inequality deliver, as in Proposition~\ref{prop:dispmon}.

\medskip
\noindent\emph{The strong-log-concavity constant in Lemma~\ref{lem:slc}.} The constant $n/(n-1)$ is
tracked as follows. On the hyperplane $H=\{\sum_i w_i=0\}$ the centred vector $w$ has $-\log$ density
$\tfrac12\|w\|^2$ (the covariance $I-\tfrac1nJ$ restricts to the identity on $H$), so the measure is
$1$-strongly-log-concave there; restricting to the convex cone $\{w_1\le\dots\le w_n\}\cap H$ adds a
convex term to $-\log$ and keeps it $1$-strongly-log-concave. Strong log-concavity is preserved under
marginalization along a linear functional \citep{prekopa1973,brascamplieb1976,saumard2014}: for an
$m$-strongly-log-concave measure and $\ell(x)=\langle v,x\rangle$, the law of $\ell$ is
$(m/\|v\|^2)$-strongly-log-concave (the marginal along the unit direction $v/\|v\|$ is $m$-slc, and
rescaling the argument by $\|v\|$ scales the constant by $1/\|v\|^2$). Here $D=w_{(k)}=\langle e_k,w\rangle$
on the cone, with $v=\mathrm{proj}_H(e_k)=e_k-\tfrac1n\mathbf 1$ and $\|v\|^2=\tfrac{n-1}n$, so
$m=1$ gives the constant $n/(n-1)$ (the cone restriction only raises it). Numerically the realized
constant is larger still, $(\log g_D)''\le-2.5$ in the cases checked, so the bound holds with room.

\bibliographystyle{plainnat}
\bibliography{../notes}

\end{document}
